\paragraph{Proposition (growing-mode obstruction).}
In the finite-\(M\) Moran MSRC model with zero background switching, the
unresolved survival of a same-background pair on a branch is governed by
\[
B_M=Q_M-\operatorname{diag}\{1/\max(M-k,1):0\le k\le M\}.
\]
For the neutral Moran rates, the interior block on
\(1\le k\le M-1\) is an irreducible tridiagonal matrix with a simple
spectrum. Therefore, for generic initial count distributions, its survival
function contains \(M-1\) independent exponential modes. Any exact
time-homogeneous linear realization of this branch-value family requires
dimension at least \(M-1\).

\emph{Proof.} While the pair is unresolved, the count process evolves with
the Moran generator and is killed at its background-specific coalescence rate,
which gives the displayed subgenerator. Restricting to interior counts removes
the absorbing count boundaries and leaves an irreducible tridiagonal matrix.
It is diagonally symmetrizable under the positive neutral birth and death
rates, and an irreducible symmetric tridiagonal matrix has simple eigenvalues.
Generic initial and terminal vectors have nonzero projection onto each
eigenmode, so the survival function contains \(M-1\) distinct exponentials.
A \(C\)-dimensional time-homogeneous linear realization has at most \(C\)
distinct exponential modes. Hence \(C\ge M-1\). The pair survival appears in
the balanced-tree early term \(1-s_Ls_R\), so a universal branch-value closure
for the arbitrary hierarchical recursion cannot remove these modes. \qed

\paragraph{Scope.}
This proposition rules out a universal finite-dimensional linear closure for
the branch-value functions required by the full arbitrary-branch recursion.
It does not claim that every specially chosen final scalar or symmetric
parameter setting has no accidental cancellation. Such cancellations are
special cases and do not supply a general formula for arbitrary initial count
distributions and branch parameters.

\paragraph{Corollary.}
The original hierarchical model does not admit a general exact unconditional
count-free recursion by ordinary moment summation, finitely many
generating-function evaluations, or strong lumping. The count-indexed operator
remains necessary for the general branch recursion.

\paragraph{Separate model.}
If the Moran count is replaced by a two-state frequency-regime CTMC with fixed
frequencies \(p_L,p_H\), the shared-population quartet law is exactly obtained
from the 38-state operator
\[
R\otimes I_{19}+\operatorname{blockdiag}(G_4(p_L)^T,G_4(p_H)^T).
\]
This statement concerns the regime model and is not an equivalence theorem
for the finite-\(M\) Moran process.
